Existing quadrupedal locomotion learning paradigms usually rely on extensive domain randomization to alleviate the sim2real gap and enhance robustness. It trains policies by introducing various disturbances into simulated environments to improve their adaptability under uncertainty. However, since optimal performance without disturbances often conflicts with the need to handle worst-case scenarios, this paradigm has a trade-off between optimality and robustness. This trade-off forces the learned policy to prioritize stability in diverse and challenging environments over efficiency and accuracy in no-disturbed ones, leading to overly conservative behaviors that sacrifice peak performance. Inspired by disturbance rejection control, we propose a two-stage framework that integrates policy learning with motion reference to mitigate this phenomenon. This framework enhances the conventional reinforcement learning (RL) approach by incorporating imagined transitions as reference inputs. The imagined transitions are derived from an optimal policy and a dynamics model operating within an idealized setting without disturbances. By providing this reference, we found that the policy could generate actions that are better aligned with the desired motion, which led to accelerated training and reduced tracking errors under external disturbances.